% preamble.tex – Gemensamma inställningar för LuaLaTeX-dokument
\ProvidesFile{preamble.tex}[2025-05-04 Gemensam preambel för LuaLaTeX]

% === Dokumentklass ===
\documentclass[12pt,a4paper]{article}

% === Språk och teckenkodning ===
\usepackage[swedish, english]{babel}
\usepackage{fontspec}
\setmainfont{TeX Gyre Termes}
\setsansfont{TeX Gyre Heros}
\setmonofont{TeX Gyre Cursor}

% === Typografi ===
\usepackage{microtype}
\usepackage{parskip}
\usepackage{setspace}
\onehalfspacing

% === Matematik ===
\usepackage{amsmath}
\usepackage{amsfonts}
\usepackage{amssymb}
\usepackage{mathtools}

% === Tabeller och figurer ===
\usepackage{booktabs}
\usepackage{array}
\usepackage{graphicx}
\graphicspath{{./bilder/}{./figurer/}}
\usepackage{caption}
\usepackage{subcaption}

% === Färger ===
\usepackage{xcolor}
\definecolor{maincolor}{RGB}{0, 102, 204}

% === Länkar och referenser ===
\usepackage[colorlinks=true, linkcolor=maincolor, citecolor=maincolor, urlcolor=maincolor]{hyperref}
\usepackage{cleveref}

% === Sidhuvud och sidfot ===
\usepackage{fancyhdr}
\pagestyle{fancy}
\fancyhf{}
\fancyhead[L]{\leftmark}
\fancyhead[R]{\thepage}
\renewcommand{\headrulewidth}{0.4pt}

% === Övrigt ===
\usepackage{enumitem}
\usepackage{tikz}
\usetikzlibrary{arrows.meta, positioning, shapes}

% === Egna kommandon ===
\newcommand{\vect}[1]{\mathbf{#1}}
\newcommand{\dd}{\mathrm{d}}
